\section{conclusion}
In this work, we introduced LightMem, a lightweight and efficient memory framework designed to address the significant overhead of memory systems for LLM agents.
Inspired by the multi-stage Atkinson-Shiffrin human memory model, LightMem's architecture effectively filters, organizes, and consolidates information. 
Our empirical evaluation demonstrates that this approach maintains strong task performance while sharply reducing computational costs. 
In the near future, we plan to accelerate LightMem’s update phase via offline pre-computed KV caches, reducing runtime overhead. We aim to integrate a lightweight knowledge graph memory for explicit multi-hop reasoning and structured retrieval. A multimodal memory extension will enable adaptation to visual, auditory, and textual inputs in embodied and real-world scenarios.

% 我们后续会进一步将lightmem框架进行扩展，包括在离线更新阶段KV Cache的预计算，使得更新过程进一步加速；为了更好地解决多跳推理难题，我们还会集成一个轻量级的knowledge graph-based memory；以及一个多模态的memory，来适配更多的多模态模型以及适应更多的场景，尤其是agent/具身智能等对多模态和momory需求更高的场景。

% \textbf{Offline Update Acceleration.}
% We plan to enhance the efficiency of LightMem’s update phase by incorporating pre-computed key–value (KV) caches.
% This optimization allows KV cache pre-computation to be performed offline, significantly accelerating memory consolidation and reducing runtime overhead during interactive sessions.

% \textbf{Knowledge Graph-based Memory.}
% To better address complex multi-hop reasoning challenges, we aim to integrate a lightweight knowledge graph-based memory module.
% This module will support explicit relational reasoning and structured information retrieval, enabling the model to perform more interpretable and compositional reasoning across interconnected knowledge entities.

% \textbf{Multimodal Memory Extension.}
% We further plan to develop a multimodal memory mechanism that allows LightMem to adapt to multimodal models and scenarios.
% Such an extension is crucial for embodied agents and real-world applications where visual, auditory, and textual information jointly contribute to memory formation and retrieval.

% \textbf{Parametric–Nonparametric Synergy.}
% We will explore the collaborative mechanisms between parametric and non-parametric memory components.
% To bridge the gap between the two paradigms, we aim to enable more flexible, synergistic knowledge utilization, combining the efficiency of parametric representations with the interpretability and adaptability of non-parametric storage.


